% Source: https://artofproblemsolving.com/community/c3h2363616_tesseractexe
% Stable identity: aops-topic-2363616-q03
Take the tesseract to be the set of points $(x_1,x_2,x_3,x_4)$ with $0\le x_i\le1$. Its vertices are the $2^4=16$ points with every coordinate equal to $0$ or $1$:
\begin{align*}
&(0,0,0,0),\ (1,0,0,0),\ (0,1,0,0),\ (0,0,1,0),\ (0,0,0,1),\ (1,1,0,0),\ (1,0,1,0),\ (1,0,0,1),\\
&(0,1,1,0),\ (0,1,0,1),\ (0,0,1,1),\ (1,1,1,0),\ (1,1,0,1),\ (1,0,1,1),\ (0,1,1,1),\ (1,1,1,1).
\end{align*}
This continues the pattern of the square (vertices $(0,0),(1,0),(0,1),(1,1)$) and the cube (the $8$ points with coordinates $0$ or $1$): each coordinate has $2$ choices, so there are $2^4=16$ vertices, as a tesseract should have. Two vertices are joined by an edge exactly when they differ in one coordinate. In the projection below a vertex $(x_1,x_2,x_3,x_4)$ is written $x_1x_2x_3x_4$:

\begin{asy}
usepackage("amsmath");
size(5cm);
defaultpen(fontsize(10pt));
pair e1=(1,0), e2=(0,1), e3=(0.35,0.42), e4=(1.25,0.55);
pair P(int a,int b,int c,int d){return a*e1+b*e2+c*e3+d*e4;}
for(int a=0;a<2;++a) for(int b=0;b<2;++b) for(int c=0;c<2;++c) for(int d=0;d<2;++d){
  if(a==0) draw(P(0,b,c,d)--P(1,b,c,d),gray);
  if(b==0) draw(P(a,0,c,d)--P(a,1,c,d),gray);
  if(c==0) draw(P(a,b,0,d)--P(a,b,1,d),gray);
  if(d==0) draw(P(a,b,c,0)--P(a,b,c,1),gray);
  dot(P(a,b,c,d));
}
label("$0000$",P(0,0,0,0),SW,fontsize(7pt)); label("$1111$",P(1,1,1,1),NE,fontsize(7pt));
label("$1000$",P(1,0,0,0),S,fontsize(7pt));
label("$0100$",P(0,1,0,0),W,fontsize(7pt));
\end{asy}
